% Clock localization subsection (Sean Ermilov).
% Meant to be \input{} into the Methodology section of the group report.
% Needs: amsmath, graphicx (both already in the template) and the entries in
% localization.bib added to references.bib. Figures are generated with
%   .venv/bin/python scripts/make_localization_figures.py
% into outputs/report_figures/ and should be copied next to the report as
% figures/localization_*.png.

\subsection{Clock localization}
\label{sec:localization}

The first stage finds the clock face so that the later stages can work
relative to its centre and radius. It returns a circle $(c_x, c_y, r)$ in the
pixel coordinates of the original photograph, or reports that no clock was
found. Every image goes through the same fixed procedure; no parameters are
chosen per image and no positions are annotated by hand. The parameters were
set by inspecting results on the training images only.

\paragraph{Resizing.}
The photographs range from about $250$ to $10000$ pixels wide. Each image is
shrunk with area interpolation so that its longest side is at most $512$
pixels, with scale factor $s = 512 / \max(w, h)$ (smaller images are not
enlarged, $s = 1$). This makes all pixel-based parameters below mean the same
thing for every image and keeps detection fast. A circle found at
$(c'_x, c'_y, r')$ in the resized image is returned as
$(c'_x / s,\; c'_y / s,\; r' / s)$ in the original image.

\paragraph{Preprocessing.}
The resized image is converted to grayscale and smoothed with a $5 \times 5$
median filter \citep{huang1979median}. The median filter removes fine texture
(wood grain, bark, printed dial details) while keeping step edges such as the
clock rim sharp. We compared this with no smoothing, a stronger median filter,
Gaussian and bilateral filtering \citep{tomasi1998bilateral}, contrast-limited
adaptive histogram equalization (CLAHE) \citep{zuiderveld1994clahe}, the HSV
saturation and value channels, and working sizes of $384$ and $768$ pixels.
None of these was clearly better on the training images, so the simplest
option was kept.

\paragraph{Circle candidates.}
Candidate circles are generated with the gradient-based circle Hough transform
\citep{duda1972hough, kimme1975circles, yuen1990circles} as implemented in
OpenCV \citep{bradski2000opencv}. Edge pixels from the Canny detector
\citep{canny1986edge} vote for centres along their gradient direction, and
radii are then estimated for the strongest centres. The radius is limited to
between $5\%$ and $60\%$ of the shorter image side, which covers both small
clocks on buildings and close-ups. Up to $60$ candidates are kept, because the
strongest Hough peak is often a circle formed by background texture rather
than the clock (Fig.~\ref{fig:loc-pipeline}d).

\paragraph{Radial edge score.}
Hough votes alone do not separate a clock rim from texture, so each candidate
is re-scored. We sample $N = 360$ points $\mathbf{p}_i$ on the circle, with
outward unit normals $\mathbf{u}_i$. Point $i$ is \emph{supported} if, within
$\pm 2$ pixels along $\mathbf{u}_i$, there is a Canny edge pixel whose image
gradient $\nabla I$ is nearly parallel to the radius:
\begin{equation}
  \left| \frac{\nabla I \cdot \mathbf{u}_i}{\lVert \nabla I \rVert} \right| \geq 0.8 .
\end{equation}
A clock rim is a closed contour whose gradients point towards or away from
the centre all the way round, which is the same property the gradient-based
Hough transform uses for voting \citep{kimme1975circles}. Texture produces
edges in random directions. The score is the fraction of supported points
minus the fraction expected by chance:
\begin{equation}
  S = \frac{1}{N} \sum_{i=1}^{N} \text{supported}_i \;-\; S_0, \qquad
  S_0 = 1 - \left(1 - \rho\,\kappa\right)^{5},
\end{equation}
where $\rho$ is the fraction of Canny edge pixels in the image,
$\kappa = 2 \arccos(0.8) / \pi \approx 0.41$ is the probability that a random
gradient direction passes the test, and $5$ is the number of pixels checked
per point. Without the correction $S_0$, any circle on a busy background would
receive a high score.

\paragraph{Selection.}
A clock face usually contains several concentric circles: the outer rim, the
inner edge of the bezel, and rings of the dial. These often score almost the
same. Among the candidates whose score is at least $80\%$ of the best score,
the one with the largest radius is selected, which favours the outer rim. If
the best score is below $0.25$, no circle is accepted.

\paragraph{Fallback for weak boundaries.}
When the selected circle scores below $0.45$, or no circle is accepted, the
rim is usually faint or low in contrast. In that case the candidate generation,
scoring and selection are repeated on a CLAHE-enhanced copy of the grayscale
image \citep{pizer1987ahe, zuiderveld1994clahe} (clip limit $2$, $8 \times 8$
tiles), with more permissive thresholds (Canny upper threshold $60$ instead of
$100$, $20$ instead of $30$ Hough votes). Of the two results, the one with the
higher score is kept (Fig.~\ref{fig:loc-fallback}). CLAHE is used only as a
fallback because it also strengthens background texture: as a first step it
produced more weak detections on the training images. The trigger is the same
score rule for every image.

\paragraph{Limitations.}
The method assumes that the clock face is bounded by a circle. Square or
rectangular clocks, clocks without a rim, and strongly tilted clocks (which
appear as ellipses) are only approximated or missed
(Fig.~\ref{fig:loc-examples}, bottom row). The score indicates weak
detections, but a high score does not guarantee that the right circle was
chosen; an inner dial ring can score as high as the rim.

\begin{figure*}[t]
    \centering
    \includegraphics[width=\linewidth]{figures/localization_pipeline.png}
    \caption{Localization steps on one training image: (a) resized input,
    (b) grayscale with median filter, (c) Canny edges, (d) Hough candidates
    coloured by radial edge score, (e) selected circle. Most candidates come
    from the bark texture and receive low scores. Photo by yaniv-knobel from
    \url{unsplash.com}.}
    \label{fig:loc-pipeline}
\end{figure*}

\begin{figure}[t]
    \centering
    \includegraphics[width=\linewidth]{figures/localization_fallback.png}
    \caption{Fallback for weak boundaries. Left: first pass. Middle:
    CLAHE-enhanced image with the fallback result. Right: result kept (higher
    score). Photos by sevenstormphotography from \url{pexels.com} (top) and
    yanghong-yu from \url{unsplash.com} (bottom).}
    \label{fig:loc-fallback}
\end{figure}

\begin{figure*}[t]
    \centering
    \includegraphics[width=\linewidth]{figures/localization_examples.png}
    \caption{Typical localization results (green) and failure cases (red):
    a square clock, a clock without a rim, a small clock on a tower where a
    larger wrong circle wins, and an inner dial ring chosen instead of the rim.
    Photos from \url{pexels.com}, \url{unsplash.com} and
    \url{commons.wikimedia.org}; photographers in the order shown:
    pixabay, ondrej-supitar, necip-duman, yash-parashar, ali-bokhari,
    will-grobbelaar, Stefi123, chirayu-trivedi, ihor-lypnytskyi, jason-tirta.}
    % TODO: check the licence of the Wikimedia photo (Stefi123) and add it,
    % as in the template's example caption.
    \label{fig:loc-examples}
\end{figure*}
